\section{Conclusion}
We compared adaptive exploration, single-round hybrid retrieval, and lexical packing for known-vulnerability localization using public descriptions and dataset-supplied entry titles. Those titles can contain target filenames, which defines the scope of the reported results. S1 shares nonblank source eligibility and retains every scheduled outcome.

Agent delivers substantially more annotated source and returns narrower intervals than RAG, but its complete-denominator file recall is lower by 2.5 percentage points. The joint analysis qualifies these outcomes: uncovered annotations with source available are predominantly entry points; most RAG hits without annotated source depend on broad intervals; and schema acceptance does not guarantee repository-valid coordinates. Agent's 31 failures contribute a negative 7.71 percentage points to the file contrast, offsetting the positive contribution from its other paired outcomes. Its higher tolerance-based line recall accompanies much greater cumulative input, with statistical support sensitive to repository dependence.

The comparison measures source acquisition, annotation coverage, interval extent, output validity, and inference expense in these configurations. Reporting S0 and S1 separately preserves the observed comparison change. The available evidence supports this joint diagnosis within the title-assisted task and the known, imbalanced sample.
